\begin{table}[!htbp]
\centering
\caption{First-stage coefficients controlling for SAE priority status}
\label{tab:scalar-school-value-iv-main-seven-eb-first-stage-prioritario}
\resizebox{\textwidth}{!}{%
\begin{tabular}{lccccccc}
\toprule
 & \multicolumn{2}{c}{Achievement} & \multicolumn{3}{c}{Postsecondary access} & \multicolumn{2}{c}{Postsecondary choice} \\
\cmidrule(lr){2-3} \cmidrule(lr){4-6} \cmidrule(lr){7-8}
 & Math & Verbal & Exam taking & Fin. aid app. & Higher-ed enrollment & High-premium inst. & High-premium field \\
\midrule
$\kappa^{EB}$ & 0.442*** & 0.434*** & 0.424*** & 0.405*** & 0.428*** & 0.268*** & 0.474*** \\
 & (0.008) & (0.007) & (0.006) & (0.006) & (0.006) & (0.009) & (0.007) \\
First-stage F & 3193.1 & 4218.9 & 5930.3 & 4226.1 & 5498.4 & 820.2 & 4278.6 \\
N & 92,361 & 93,980 & 135,844 & 135,844 & 135,844 & 94,487 & 94,487 \\
\bottomrule
\end{tabular}
}
\par\medskip
\footnotesize
\begin{minipage}{\textwidth}
Notes: This robustness specification adds the SAE low-income priority indicator, \textit{prioritario}, to the controls used in the corresponding main specification. It also controls for the DA-probability expected value of the same EB measure, cohort, grade-4 SIMCE math and language, gender, and age. The sample is restricted to students with observed priority status. Heteroskedasticity-robust standard errors are reported in parentheses. $^{*}p<0.10$, $^{**}p<0.05$, $^{***}p<0.01$.
\end{minipage}
\end{table}
